% =========================================================
% Density
% =========================================================

\subsection{Density}




\begin{tcolorbox}[colback=propbox, colframe=propborder, arc=2pt,
  left=6pt, right=6pt, top=4pt, bottom=4pt,
  title={\small\textbf{Definition (Dense Subset)}},
  fonttitle=\small\bfseries]
\begin{definition}
\label{def:dense-subset}
Let $D \subseteq S \subseteq \mathbb{R}$. The set $D$ is said to be
\emph{dense in $S$} if for every $x \in S$ and every $\varepsilon > 0$
there exists $d \in D$ such that
\[
|x-d|<\varepsilon.
\]
\end{definition}
\end{tcolorbox}

\begin{remark*}[Standard quantified statement]
\[
D \text{ is dense in } S
\quad\Longleftrightarrow\quad
\forall x \in S \; \forall \varepsilon > 0 \; \exists d \in D \;
\bigl(|x-d|<\varepsilon\bigr).
\]
\end{remark*}

\begin{remark*}[Definition predicate reading]
\[
\operatorname{DenseSubset}(D,S).
\]
Equivalently,
\[
\forall x \in S \; \forall \varepsilon > 0 \; \exists d \in D \;
\bigl(d(x,d)<\varepsilon\bigr).
\]
\end{remark*}

\begin{remark*}[Negated quantified statement]
\[
D \text{ is not dense in } S
\quad\Longleftrightarrow\quad
\exists x \in S \; \exists \varepsilon_0 > 0 \; \forall d \in D \;
\bigl(|x-d| \ge \varepsilon_0\bigr).
\]
\end{remark*}

\begin{remark*}[Failure mode decomposition]
\[
D \text{ is not dense in } S
\quad\Longleftrightarrow\quad
\underbrace{\exists x \in S \; \exists \varepsilon_0 > 0 \; \forall d \in D \;
\bigl(|x-d| \ge \varepsilon_0\bigr)}_{\text{some point of $S$ has an open neighborhood missing $D$ entirely}}.
\]
\end{remark*}

\begin{remark*}[Negation predicate reading]
\[
\neg \operatorname{DenseSubset}(D,S)
\quad\Longleftrightarrow\quad
\exists x \in S \; \exists \varepsilon_0 > 0 \; \forall d \in D \;
\bigl(d(x,d)\ge \varepsilon_0\bigr).
\]
\end{remark*}

\begin{remark*}[Failure modes]\hfill
\begin{itemize}
  \item \textbf{Failure of approximation.} There exists a point of $S$ that cannot~\\ be approximated arbitrarily closely by points of $D$.
  \item \textbf{Failure of local presence.} There exists a positive radius around some point of $S$ whose entire neighborhood misses $D$.
\end{itemize}
\end{remark*}






\input{volume-iii/analysis/bounding/notes/completeness/figure-dense-subset}





\begin{remark*}[Interpretation]
Density is an approximation property inside $S$. The definition does not require that a point
of $S$ already lie in $D$; it requires only that every point of $S$ can be approached
arbitrarily closely by points of $D$.

Equivalently, no matter how small a positive radius is chosen around a point of $S$, that
neighborhood contains at least one point of $D$. A dense subset may therefore be much smaller
than $S$ while still leaving no open gap inside it.
\end{remark*}



\begin{remark*}[Application of Dense Subset]
This definition formalizes the concept that the subset $D$ is "everywhere present" within the superset $S$. While $D$ might have gaps (such as $\mathbb{Q}$ within $\mathbb{R}$), these gaps are infinitesimally small.

\smallskip
\noindent
The quantified statement
\[
\forall x \in S, \forall \varepsilon > 0, \exists d \in D \; \bigl( |x - d| < \varepsilon \bigr)
\]
means that no element of $S$ is "isolated" from $D$. No matter how tightly you draw a tolerance boundary ($\varepsilon$) around any point $x \in S$, you are guaranteed to find a representative from $D$ lurking inside that boundary.

\smallskip
\noindent
This concept is mathematically identical to stating that the closure of $D$ (relative to $S$) is $S$ itself: $\operatorname{Cl}_S(D) = S$. In real analysis, this allows us to approximate any real number using only rational numbers, which is fundamental for constructing the reals via Cauchy sequences.
\end{remark*}







\begin{tcolorbox}[colback=propbox, colframe=propborder, arc=2pt,
  left=6pt, right=6pt, top=4pt, bottom=4pt,
  title={\small\textbf{Definition (Order Dense Subset)}},
  fonttitle=\small\bfseries]
\begin{definition}
\label{def:order-dense-subset}
Let $A \subseteq X$, where $X$ is a linearly ordered set. The subset $A$ is said to be
\emph{order dense in $X$} if every nonempty open interval in $X$ contains a point of $A$.
\end{definition}
\end{tcolorbox}

\begin{remark*}[Standard quantified statement]
\[
A \text{ is order dense in } X
\quad\Longleftrightarrow\quad
\forall a,b \in X\;
\bigl(
a<b \rightarrow \exists c \in A\; (a<c<b)
\bigr).
\]
\end{remark*}

\begin{remark*}[Negated quantified statement]
\[
A \text{ is not order dense in } X
\quad\Longleftrightarrow\quad
\exists a,b \in X\;
\bigl(
a<b \land \forall c \in A\; (c \le a \lor b \le c)
\bigr).
\]
\end{remark*}

\begin{remark*}[Failure mode decomposition]
\[
A \text{ is not order dense in } X
\quad\Longleftrightarrow\quad
\underbrace{\exists a,b \in X\;
\bigl(
a<b \land \forall c \in A\; (c \le a \lor b \le c)
\bigr)}_{\text{there is a nonempty open interval in $X$ containing no point of $A$}}.
\]
\end{remark*}

\begin{remark*}[Failure modes]\hfill
\begin{itemize}
  \item \textbf{Failure of interval penetration.} There exists a pair of points $a<b$ in $X$ such that no element of $A$ lies strictly between them.
  \item \textbf{Failure of order presence.} There exists a nonempty open interval in $X$ that is missed entirely by $A$.
\end{itemize}
\end{remark*}

\begin{remark*}[Interpretation]
Order density says that the subset $A$ leaves no interval gap in $X$. Whenever two elements
of $X$ satisfy $a<b$, there is an element of $A$ strictly between them.

The emphasis is on order, not on distance. The definition does not say that points of $A$ lie
arbitrarily close to every point in a metric sense; it says that no nonempty open interval of
$X$ is missed entirely by $A$.
\end{remark*}

















\begin{theorem}[Density of the Rational Numbers in $\mathbb{R}$]
\label{thm:density-of-rationals-in-reals}
For every $a,b \in \mathbb{R}$ with $a<b$, there exists $q \in \mathbb{Q}$ such that
\[
a<q<b.
\]
\end{theorem}

\begin{remark*}[Standard quantified statement]
\[
\forall a,b \in \mathbb{R}\;
\bigl(
a<b \rightarrow \exists q \in \mathbb{Q}\; (a<q<b)
\bigr).
\]
\end{remark*}

\begin{remark*}[Theorem predicate reading]
\[
\underbrace{\operatorname{DenseSubset}(\mathbb{Q},\mathbb{R})}_{\text{the rational numbers are dense in the real line}}
\quad\Longleftrightarrow\quad
\underbrace{\forall a,b \in \mathbb{R}\;
\bigl(
a<b \rightarrow \exists q \in \mathbb{Q}\; (a<q<b)
\bigr)}_{\text{every nonempty open interval in $\mathbb{R}$ contains a rational number}}.
\]
\end{remark*}

\begin{remark*}[Negated quantified statement]
\[
\exists a,b \in \mathbb{R}\;
\bigl(
a<b \land \forall q \in \mathbb{Q}\; (q \le a \lor b \le q)
\bigr).
\]
\end{remark*}

\begin{remark*}[Failure mode decomposition]
\[
\exists a,b \in \mathbb{R}\;
\underbrace{\bigl(
a<b \land \forall q \in \mathbb{Q}\; (q \le a \lor b \le q)
\bigr)}_{\text{there is a nonempty open interval in $\mathbb{R}$ containing no rational number}}.
\]
\end{remark*}

\begin{remark*}[Negation predicate reading]
\[
\neg \operatorname{DenseSubset}(\mathbb{Q},\mathbb{R})
\quad\Longleftrightarrow\quad
\exists a,b \in \mathbb{R}\;
\bigl(
a<b \land \forall q \in \mathbb{Q}\; (q \le a \lor b \le q)
\bigr).
\]
\end{remark*}

\begin{remark*}[Failure modes]\hfill
\begin{itemize}
  \item \textbf{Failure of interval penetration.} There exists a pair of real numbers $a<b$ such that no rational number lies strictly between them.
  \item \textbf{Failure of density.} There exists a nonempty open interval in $\mathbb{R}$ that is missed entirely by $\mathbb{Q}$.
\end{itemize}
\end{remark*}

\begin{remark*}[Interpretation]
This theorem says that rational numbers are spread throughout the real line with no interval
gap. Whenever two distinct real numbers are given, at least one rational number lies between
them.

The point is not that $\mathbb{Q}=\mathbb{R}$, but that $\mathbb{Q}$ is present everywhere in
$\mathbb{R}$ at arbitrarily small interval scales. Every nonempty open interval in $\mathbb{R}$
contains rational points.
\end{remark*}







\begin{theorem}[Density of the Irrational Numbers in $\mathbb{R}$]
\label{thm:density-of-irrationals-in-reals}
For every $a,b \in \mathbb{R}$ with $a<b$, there exists
$s \in \mathbb{R}\setminus\mathbb{Q}$ such that
\[
a<s<b.
\]
\end{theorem}

\begin{remark*}[Standard quantified statement]
\[
\forall a,b \in \mathbb{R}\;
\bigl(
a<b \rightarrow \exists s \in \mathbb{R}\setminus\mathbb{Q}\; (a<s<b)
\bigr).
\]
\end{remark*}

\begin{remark*}[Theorem predicate reading]
\[
\begin{aligned}
&\underbrace{\operatorname{DenseSubset}(\mathbb{R}\setminus\mathbb{Q},\mathbb{R})}
_{\text{the irrational numbers are dense in the real line}}
&\Longleftrightarrow
\\[0.5ex]
&\underbrace{
\forall a,b \in \mathbb{R}\;
\bigl(
a<b \rightarrow \exists s \in \mathbb{R}\setminus\mathbb{Q}\; (a<s<b)
\bigr)
}_{\text{every nonempty open interval in $\mathbb{R}$ contains an irrational number}}.
\end{aligned}
\]
\end{remark*}

\begin{remark*}[Negated quantified statement]
\[
\exists a,b \in \mathbb{R}\;
\bigl(
a<b \land \forall s \in \mathbb{R}\setminus\mathbb{Q}\; (s \le a \lor b \le s)
\bigr).
\]
\end{remark*}

\begin{remark*}[Failure mode decomposition]
\[
\exists a,b \in \mathbb{R}\;
\underbrace{\bigl(
a<b \land \forall s \in \mathbb{R}\setminus\mathbb{Q}\; (s \le a \lor b \le s)
\bigr)}_{\text{there is a nonempty open interval in $\mathbb{R}$ containing no irrational number}}.
\]
\end{remark*}

\begin{remark*}[Negation predicate reading]
\[
\neg \operatorname{DenseSubset}(\mathbb{R}\setminus\mathbb{Q},\mathbb{R})
\quad\Longleftrightarrow\quad
\exists a,b \in \mathbb{R}\;
\bigl(
a<b \land \forall s \in \mathbb{R}\setminus\mathbb{Q}\; (s \le a \lor b \le s)
\bigr).
\]
\end{remark*}

\begin{remark*}[Failure modes]\hfill
\begin{itemize}
  \item \textbf{Failure of interval penetration.} There exists a pair of real numbers $a<b$ such that no irrational number lies strictly between them.
  \item \textbf{Failure of density.} There exists a nonempty open interval in $\mathbb{R}$ that is missed entirely by $\mathbb{R}\setminus\mathbb{Q}$.
\end{itemize}
\end{remark*}

\begin{remark*}[Interpretation]
This theorem says that irrational numbers are spread throughout the real line with no interval
gap. Whenever two distinct real numbers are given, at least one irrational number lies between
them.

The point is not merely that irrational numbers exist, but that they occur everywhere in
$\mathbb{R}$ at arbitrarily small interval scales. Every nonempty open interval in $\mathbb{R}$
contains irrational points.
\end{remark*}




\begin{corollary}
\label{cor:every-open-interval-contains-rational-and-irrational}
Every open interval in $\mathbb{R}$ contains both a rational number and an irrational number.
\end{corollary}

\begin{remark*}[Standard quantified statement]
\[
\forall a,b \in \mathbb{R}\;
\Bigl[
a<b \rightarrow
\bigl(
\exists q \in \mathbb{Q}\; (a<q<b)
\land
\exists s \in \mathbb{R}\setminus\mathbb{Q}\; (a<s<b)
\bigr)
\Bigr].
\]
\end{remark*}

\begin{remark*}[Theorem predicate reading]
\[
\underbrace{\operatorname{DenseSubset}(\mathbb{Q},\mathbb{R})}_{\text{the rational numbers are dense in }\mathbb{R}}
\land
\underbrace{\operatorname{DenseSubset}(\mathbb{R}\setminus\mathbb{Q},\mathbb{R})}_{\text{the irrational numbers are dense in }\mathbb{R}}.
\]
Equivalently,
\[
\underbrace{\forall a,b \in \mathbb{R}\;
\Bigl[
a<b \rightarrow
\bigl(
\exists q \in \mathbb{Q}\; (a<q<b)
\land
\exists s \in \mathbb{R}\setminus\mathbb{Q}\; (a<s<b)
\bigr)
\Bigr]}_{\text{every nonempty open interval in $\mathbb{R}$ contains both kinds of points}}.
\]
\end{remark*}

\begin{remark*}[Negated quantified statement]
\[
\exists a,b \in \mathbb{R}\;
\Bigl[
a<b \land
\Bigl(
\forall q \in \mathbb{Q}\; (q \le a \lor b \le q)
\;\lor\;
\forall s \in \mathbb{R}\setminus\mathbb{Q}\; (s \le a \lor b \le s)
\Bigr)
\Bigr].
\]
\end{remark*}

\begin{remark*}[Failure mode decomposition]
\[
\exists a,b \in \mathbb{R}\;
\Bigl[
a<b \land
\Bigl(
\underbrace{\forall q \in \mathbb{Q}\; (q \le a \lor b \le q)}_{\text{the interval contains no rational number}}
\;\lor\;
\underbrace{\forall s \in \mathbb{R}\setminus\mathbb{Q}\; (s \le a \lor b \le s)}_{\text{the interval contains no irrational number}}
\Bigr)
\Bigr].
\]
\end{remark*}

\begin{remark*}[Negation predicate reading]
\[
\neg\Bigl(
\operatorname{DenseSubset}(\mathbb{Q},\mathbb{R})
\land
\operatorname{DenseSubset}(\mathbb{R}\setminus\mathbb{Q},\mathbb{R})
\Bigr)
\]
\[
\Longleftrightarrow\quad
\exists a,b \in \mathbb{R}\;
\Bigl[
a<b \land
\Bigl(
\forall q \in \mathbb{Q}\; (q \le a \lor b \le q)
\;\lor\;
\forall s \in \mathbb{R}\setminus\mathbb{Q}\; (s \le a \lor b \le s)
\Bigr)
\Bigr].
\]
\end{remark*}

\begin{remark*}[Failure modes]\hfill
\begin{itemize}
  \item \textbf{Failure of rational presence.} There exists a nonempty open interval in $\mathbb{R}$ containing no rational number.
  \item \textbf{Failure of irrational presence.} There exists a nonempty open interval in $\mathbb{R}$ containing no irrational number.
\end{itemize}
\end{remark*}

\begin{remark*}[Interpretation]
This corollary combines the two density theorems into a single local statement about the real
line. No nonempty open interval contains only rational points, and no nonempty open interval
contains only irrational points.

The point is that both kinds of numbers occur everywhere in $\mathbb{R}$. Every open interval,
no matter how small, contains representatives of both $\mathbb{Q}$ and
$\mathbb{R}\setminus\mathbb{Q}$.
\end{remark*}




\input{volume-iii/analysis/bounding/notes/completeness/figure-zoom-invariance}


\begin{remark*}[Zoom invariance and structural consequences]~\\
Density in $\mathbb{R}$ has a scale-invariant consequence: no matter how small a nonempty
open interval becomes, it still contains both rational and irrational points. The interleaving
of $\mathbb{Q}$ and $\mathbb{R}\setminus\mathbb{Q}$ is therefore not a large-scale phenomenon.
It persists at every interval scale.

This has three important consequences. First, the continuum does not become locally discrete
under magnification. For every nonempty open interval $(a,b)$, both
\[
\mathbb{Q}\cap(a,b)
\qquad\text{and}\qquad
(\mathbb{R}\setminus\mathbb{Q})\cap(a,b)
\]
remain nonempty, and in fact infinite. Density is therefore stable under repeated restriction
to smaller intervals.

Second, this behavior sharply distinguishes the real line from discrete numerical models.
Floating-point systems form a discrete set of representable values, so repeated magnification
eventually reaches a scale at which adjacent representable numbers leave no further point
strictly between them. At that stage, interval penetration fails in the discrete model even
though it continues to hold in $\mathbb{R}$.

Third, the Archimedean Property supplies the scale mechanism behind this behavior. Given any
positive scale $\varepsilon$, there exists $n\in\mathbb{N}$ with
\[
\frac{1}{n}<\varepsilon.
\]
This makes it possible to construct rational approximations at arbitrarily small interval
scales, which is one of the basic engines behind density arguments.

The resulting contrast is structural. Smooth mathematical models are built in a setting where
interval penetration persists at every scale, whereas finite computational models eventually
encounter a smallest representable separation. The transition from continuum behavior to
discrete behavior occurs exactly where nonempty intervals in the numerical model cease to
contain further \\representable points.
\end{remark*}

\input{volume-iii/analysis/bounding/notes/completeness/figure-density-chain}